\chapter{Limitations}
\label{ch:limitations}

This chapter states explicitly what the formalization does \emph{not}
provide.

\section{Operational Runtimes}

The formalization provides mathematical specifications and correctness
theorems.  It does not provide an executable PLN inference engine,
performance benchmarks on large knowledge bases, or integration with
MeTTa or Hyperon runtimes.

These are operational concerns that depend on runtime infrastructure
not yet available.  The theorem-level results constrain what any
correct implementation must satisfy, but they do not constitute an
implementation.  The gap between a certified specification and a
deployable system is real, and we do not pretend otherwise.

\section{Finite-Domain Restriction}

The evidence carriers, view projections, and inference-control
theorems are formalized over finite domains.  An infinitary layer
exists in the codebase but has not been brought to chapter-level
parity.  Measure-theoretic extensions---particularly for
continuous-valued evidence and distributional inference over
non-conjugate families---require additional infrastructure that is
within reach of the existing kernel design but is not yet realized
(see Chapter~\ref{ch:future}, \S\ref{sec:future-measure}).

This restriction is not merely a matter of generality.  Several
important applications---continuous sensor fusion, Gaussian process
priors, density estimation---require the infinitary extension to be
fully rigorous.  Until it is completed, claims about these domains
must be understood as extrapolations from the finite theory.

\section{Empirical Scope}

The benchmark validation chapters
(Chapters~\ref{ch:carrier-validation}--\ref{ch:pplbench}) cover a
specific set of domains: forecasting, sensor drift, fault diagnosis,
anomaly detection, and Bayesian-network inference.  However, several
important application domains remain unvalidated:

\begin{itemize}
\item Drug-repurposing knowledge-graph benchmarks
  (Hetionet/Rephetio)
\item Genomic regulatory pipelines
\item Large-scale knowledge-base reasoning
\item Real-time streaming inference with concept drift
\end{itemize}

\noindent
These gaps do not indicate theoretical limitations but reflect the
practical scope of empirical validation completed to date.  The
formalization provides the theoretical foundation for these
applications, but the experimental demonstration is incomplete.

\section{Universal Prediction}

The Universal Prediction module explores connections between the
world-model calculus and algorithmic information theory.  This module
is work-in-progress and is not part of the stable core narrative.
Its theorems are not referenced by the main development, and its
sorry count is nonzero.

\section{Codebase Organization}

The Lean codebase reflects the history of the development rather than
an idealized dependency graph.  A comprehensive refactoring into a
``mathlibified'' library structure---with uniform naming conventions,
minimal imports, and docstring coverage---is a worthy future goal but
is explicitly not attempted here.
